%%%%%%%%%%%%%%%%%%%%%%%%%%%%%%%%%%%%%%%%%%%%%%%%%%%%%%%%%%%%%%%%%%%%%%%%%%%%%%%%
%2345678901234567890123456789012345678901234567890123456789012345678901234567890
%        1         2         3         4         5         6         7         8

\documentclass[letterpaper, 10 pt, conference]{ieeeconf}  % Comment this line out if you need a4paper

%\documentclass[a4paper, 10pt, conference]{ieeeconf}      % Use this line for a4 paper

\IEEEoverridecommandlockouts                              % This command is only needed if 
                                                          % you want to use the \thanks command

\overrideIEEEmargins                                      % Needed to meet printer requirements.

%In case you encounter the following error:
%Error 1010 The PDF file may be corrupt (unable to open PDF file) OR
%Error 1000 An error occurred while parsing a contents stream. Unable to analyze the PDF file.
%This is a known problem with pdfLaTeX conversion filter. The file cannot be opened with acrobat reader
%Please use one of the alternatives below to circumvent this error by uncommenting one or the other
%\pdfobjcompresslevel=0
%\pdfminorversion=4

% See the \addtolength command later in the file to balance the column lengths
% on the last page of the document

% The following packages can be found on http:\\www.ctan.org
\usepackage{graphicx} % for pdf, bitmapped graphics files
%\usepackage{epsfig} % for postscript graphics files
%\usepackage{mathptmx} % assumes new font selection scheme installed
%\usepackage{times} % assumes new font selection scheme installed
\usepackage{amsmath} % assumes amsmath package installed
\usepackage{multirow}
%\usepackage{amssymb}  % assumes amsmath package installed
\makeatletter
\let\NAT@parse\undefined
\makeatother
\usepackage[backref]{hyperref}

\title{\LARGE \bf
From Balancing to Tracking: Predictive Control with Structured Residual Dynamics for Indirect Drone-Ball Manipulation
}


\author{Mingjiang Liu$^{1}$, Xiaozhen Zhang$^{1}$, Stephen L. Smith$^{2}$, and Hailong Huang$^{1}$% <-this % stops a space
% \thanks{*This work was supported by the Research Centre for Low Altitude Economy (RCLAE), The Hong Kong Polytechnic University.}% <-this % stops a space
\thanks{$^{1}$Mingjiang Liu, Xiaozhen Zhang, and Hailong Huang are with the Department of Aeronautical and Aviation Engineering, The Hong Kong Polytechnic University, Hong Kong, China}
\thanks{$^{2}$Stephen L. Smith is with the Department of Electrical and Computer Engineering, University of Waterloo, Waterloo, Canada}
}


\begin{document}



\maketitle
\thispagestyle{empty}
\pagestyle{empty}


%==============================================================================
\begin{abstract}
Indirect drone-ball manipulation is sensitive to actuation delay and model mismatch because the ball moves
freely along the beam. This paper presents a unified predictive-control framework for target-position balancing
and trajectory tracking by treating balancing as constant-reference tracking. The
augmented nominal model combines ball-beam dynamics with an identified first-order-plus-dead-time low-level
actuation response, while a state predictor handles pending
commands. To address real-system mismatch, we learn a bounded and smooth ball-acceleration
residual composed of a position-dependent beam-profile correction and a friction-inspired dissipative term. The
same frozen residual is used in both the state predictor and nonlinear model predictive controller (NMPC). In
ideal simulation, nominal NMPC yields lower errors than a nonlinear feedforward--feedback (NFFB) baseline for
every tested task. On held-out real-world data, the structured residual reduces position and velocity root mean
square error (RMSE) by
$19.7\%$ and $20.8\%$ over a one-second causal rollout. In physical experiments, residual-aware NMPC reduces
overall steady-state error by $55.2\%$ for constant targets and overall mean absolute error (MAE) by $31.1\%$ for
dynamic references, while raising response gain from $0.855$ to $0.950$. These results show that
structured residual dynamics improve control-relevant state prediction and closed-loop accuracy on the tested
system.



\end{abstract}

% Reserve the next column top (the first-page right column) for the task illustration.
\suppressfloats[t]
\begin{figure}[t]
    \centering
    \includegraphics[width=\columnwidth]{figures/task_illustration_paper.pdf}
    \caption{Illustration of the indirect drone-ball trajectory tracking setup. By moving vertically, the tethered drone changes
    the beam inclination to guide the freely rolling ball along a desired position trajectory
    (purple dashed line). Snapshots at $t_1$--$t_4$ illustrate the system's temporal evolution,
    with earlier states shown more transparently.}
    \label{fig:task_illustration}
\end{figure}


%==============================================================================
\section{INTRODUCTION}
Beyond mobility and sensing, aerial robots are increasingly expected to physically interact with and alter 
their environments through manipulation. Aerial manipulation is therefore a key  direction for broadening 
operational scope and practical utility of aerial robots~\cite{ollero2022past}. Recent work has advanced direct
aerial manipulation in precision, versatility, and autonomy. Representative systems have demonstrated
millimeter-level pick and peg-in-hole operations~\cite{wang2024millimeter}, enabled writing and pick-and-place
through a common end-effector-centric framework~\cite{he2025flyingHand}, and integrated flight and manipulation
for autonomous grasping, perching, door opening, and transport~\cite{wu2026handLike}. Despite these advances,
these systems largely follow a direct-interaction paradigm: an onboard arm or gripper acts on an object that is
then rigidly coupled to the end effector or engaged with a largely stationary or quasi-static contact target.

Tool use in nature illustrates how acting through an intermediate object can extend an animal's
physical capabilities. New Caledonian crows manufacture hooked tools to extract prey~\cite{hunt1996manufacture},
have been observed bending wire into hooks to retrieve food containers~\cite{weir2002shaping},
and can use tools sequentially to obtain other tools before reaching food~\cite{wimpenny2009cognitive}.
These behaviors motivate a complementary direction for aerial robotics: extending physical interaction
through tools and intermediate mechanisms. Beyond selecting, constructing, and sequencing tools,
a central control problem is how to use a given tool to regulate object motion. This becomes
indirect dynamic aerial manipulation when aerial actuation drives an intermediate mechanism
while the object retains its own dynamics. Aerial-Balance-Bench provides a controlled and reproducible
testbed for this execution problem with a fixed tool~\cite{liu2026aerialBalanceBench}: a vertically
constrained tethered drone tilts a beam to guide a freely rolling ball toward a target position
or along a reference trajectory, as illustrated in Fig.~\ref{fig:task_illustration}.
The ball is neither grasped by nor rigidly attached to the drone;
its motion is controlled through the beam. Although simplified relative to unconstrained flight,
the setup retains the coupling among aerial actuation, tool motion, and object dynamics.

Aerial-Balance-Bench also leaves two control questions open. First, although it defined two tasks, its
reference controllers were evaluated only on target-position balancing in simulation and on the real system.
It therefore did not show whether one controller can cover static regulation and dynamic trajectory tracking
under a common formulation. Second, its delay-aware state predictor compensated for the identified command
latency, but control remained based on nominal ball-beam dynamics. The physical experiments showed lower
performance at off-center targets, indicating that unmodeled or nonideal dynamics continued to affect the
closed loop~\cite{liu2026aerialBalanceBench}. For dynamic references, such mismatch can accumulate over a
prediction horizon and appear as tracking error, steady offset, or reduced response amplitude.

These observations lead to the central question of this work: can one predictive controller handle both constant
and time-varying references while reducing the effects of real-system model mismatch? To answer this question,
we first express target-position balancing as the constant-reference case of a unified trajectory-tracking problem.
We then develop a nonlinear model predictive controller (NMPC) based on augmented nominal dynamics that include
the identified first-order-plus-dead-time (FOPDT) low-level actuation response. The first-order velocity response
is introduced as an augmented model state, whereas the dead time is handled by propagating the measured state
through the pending command history. NMPC starts from this delay-compensated predicted state and uses reference
preview over its horizon.

Even with the augmented nominal model, some real ball-beam effects remain unmodeled. We therefore construct a
structured residual for the ball acceleration from real-system data. It combines a ball-position-dependent beam
profile correction with a friction-inspired dissipative residual. Both terms are kept bounded and smooth so that
the residual can be evaluated and differentiated throughout real-time prediction. Crucially, the same residual
model is integrated into both the delay state predictor and the NMPC prediction model. This defines the
residual-aware predictive controller and keeps the short delay-compensation rollout consistent with the following
NMPC horizon. Comparing it with nominal predictive control tests whether better control-related dynamics
prediction leads to more accurate closed-loop balancing and tracking in this system.

The main contributions of this work are summarized as three folds:
\begin{itemize}
    \item We formulate target-position balancing and dynamic trajectory tracking as one reference-tracking
    problem and develop a unified NMPC based on augmented nominal dynamics, including an identified FOPDT
    low-level actuation response.
    \item We construct a bounded and smooth structured residual dynamics model for the drone-ball system and
    integrate it consistently into the delay state predictor and NMPC, yielding residual-aware predictive control.
    \item We conduct a three-part evaluation showing that the unified NMPC outperforms NFFB in ideal simulation,
    the structured residual improves causal multi-step prediction on real data, and residual-aware NMPC improves
    real-world balancing and trajectory tracking over nominal NMPC while retaining real-time feasibility.
\end{itemize}

The rest of the paper is organized as follows. Section~\ref{sec:related_work} reviews related work.
Section~\ref{sec:problem_formulation} presents the unified task formulation and augmented nominal dynamics.
Section~\ref{sec:residual_aware_control} introduces the structured residual dynamics and their integration into
predictive control. Section~\ref{sec:experiments} evaluates unified tracking in simulation, offline prediction
on real data, and real-world closed-loop performance. Section~\ref{sec:conclusions} concludes the paper.


%==============================================================================
\section{Related Work}
\label{sec:related_work}
\subsection{Aerial Manipulation}
Aerial manipulation has primarily relied on onboard arms or grippers for direct object interaction. A hardware
benchmark compared representative platforms in accuracy, execution time, manipulation capability, and impact
response~\cite{2020_ral_benchmarks_for_aerial_manipulation}. Recent advances include predictive optimization for
millimeter-level pick and peg-in-hole operations~\cite{wang2024millimeter}; robust control and whole-body planning
for omnidirectional manipulation at arbitrary base orientations~\cite{2025_ijrr_omnidirectional_aerial_manipulator};
the end-effector-centric Flying Hand framework for teleoperated and learned tasks~\cite{he2025flyingHand}; onboard
perception and soft gripping for high-speed grasping of static and moving targets~\cite{2024_npj_robotics_aerial_grasp};
and HI-ARM's integrated flight and hand-like grasping for autonomous multi-task
interaction~\cite{wu2026handLike}. Cooperative direct manipulation has also progressed: FlyingToolbox enables
vertically stacked MAVs to exchange tools with sub-centimeter docking accuracy despite strong
downwash~\cite{cao2025proximal}. Flexible connections provide another approach. RotorTM models the full-system and
slack-taut hybrid dynamics of cable-suspended loads~\cite{2023_tro_rotortm}. Impact-aware planning and hybrid NMPC
handle cable-mode transitions~\cite{2024_tro_impact_aware_payloads}. Multi-robot cable suspension supports
six-degree-of-freedom physical human-robot collaborative manipulation~\cite{2024_tro_human_aware_payload_transportation},
while online whole-body kinodynamic planning and tension-aware control enable agile cooperative load manipulation
through narrow passages~\cite{sun2025agile}. At the benchmark level, AIR-VLA addresses long-horizon
vision-language-action tasks for arm-equipped aerial systems~\cite{2026_arxiv_air_vla}, whereas
Aerial-Balance-Bench targets indirect dynamic manipulation~\cite{liu2026aerialBalanceBench}. Unlike most settings
above, where the object is grasped or physically connected to the robot, Aerial-Balance-Bench considers a ball that moves freely
along a beam and is controlled indirectly by the drone, broadening the task scope of aerial manipulation. However,
its reference controllers neither unified balancing and trajectory tracking within one framework nor addressed
real-system model mismatch. This work targets these two gaps.

\subsection{Learning-Augmented Dynamics for Control}
Learning-augmented control compensates for effects that are difficult to model from first principles. Neural
Lander augments nominal dynamics with a neural residual for near-ground aerodynamic effects~\cite{shi2019neural}.
Neural-Swarm2 learns inter-vehicle aerodynamic interactions for swarm planning and
control~\cite{shi2021neural}. Data-driven MPC uses Gaussian-process residuals for high-speed quadrotor
tracking~\cite{torrente2021data}. Real-time Neural MPC embeds high-capacity neural dynamics models in online
optimization~\cite{salzmann2023real}. Other approaches adapt during deployment. Neural-Fly learns a shared
aerodynamic representation offline and updates low-dimensional coefficients online for changing
winds~\cite{o2022neural}. Uncertainty-aware MPC combines offline and online dynamics learning while selecting
actions that improve both control and data collection~\cite{saviolo2023active}. Simultaneous system identification
and MPC learns unknown dynamics from closed-loop data with dynamic-regret guarantees~\cite{zhou2025simultaneous}.
These studies primarily learn vehicle-side aerodynamic effects, whereas we target mismatch in the
manipulated-object dynamics. We construct a control-oriented residual that encodes task-relevant physical
structure and use it consistently in both delay-aware state prediction and NMPC.


%==============================================================================
\section{Problem Formulation}
\label{sec:problem_formulation}
\subsection{Unified Task Formulation}
The system considered in this work follows the physical configuration of
Aerial-Balance-Bench~\cite{liu2026aerialBalanceBench}. For the tasks considered here, a rope connects the drone to
one end of a beam whose opposite end remains stationary, and the drone is constrained to vertical translation.
The resulting vertical motion changes the beam inclination and thereby regulates the
motion of a ball along the beam. Let $p_b(t)$ denote
the ball position measured along the beam and let $\theta(t)$ denote the beam inclination. We define the feasible
ball-motion interval and the reference interval as
\begin{equation}
    \mathcal{P}=[p_{\min},p_{\max}], \qquad
    \mathcal{P}_g\subset\mathcal{P},
    \label{eq:feasible_ball_interval}
\end{equation}
where $p_b(t)\in\mathcal{P}$ during normal task execution and $\mathcal{P}_g$ leaves a safety margin from the beam
ends. This constrained setup does not represent arbitrary three-dimensional flight, but it retains the main
coupling among drone vertical motion, beam inclination, and ball motion.

Aerial-Balance-Bench defines two task families: target-position balancing and trajectory tracking. We represent both 
with a common desired ball-position signal \(p_g(t)\). For an evaluation interval $t\in[0,T]$, the reference is assumed 
to be available at the current time and over the required prediction horizon. We define the unified reference
$p_g\colon[0,T]\rightarrow\mathcal{P}_g$ as
\begin{equation}
    p_g(t)=
    \begin{cases}
        p_g^\star, & \text{target-position balancing},\\
        p_g^{\mathrm{tr}}(t), & \text{trajectory tracking}.
    \end{cases}
    \label{eq:unified_position_reference}
\end{equation}
where $p_g^\star$ is a constant target and $p_g^{\mathrm{tr}}(t)$ is a bounded time-varying reference. The same
definition therefore covers static regulation and dynamic tracking without changing the task variables.

The common position-tracking error is $e_p(t)$. We measure task performance by the normalized
position error over the evaluation interval and require the ball to remain within the feasible beam region:
\begin{equation}
\begin{aligned}
    e_p(t) &= p_b(t)-p_g(t),\\
    J_{\mathrm{task}}
    &=\frac{1}{T}\int_{0}^{T}|e_p(t)|\,\mathrm{d}t,
\end{aligned}
    \label{eq:unified_task_objective}
\end{equation}
with $p_b(t)\in\mathcal{P}$ throughout the evaluation. For a constant reference, the ideal regulation condition
is
\begin{equation}
    p_g(t)\equiv p_g^\star
    \quad\Longrightarrow\quad
    \lim_{t\rightarrow\infty}|p_b(t)-p_g^\star|=0.
    \label{eq:balancing_special_case}
\end{equation}
Thus, target-position balancing is the constant-reference case of the unified task, whereas a nonconstant
$p_g(t)$ gives trajectory tracking under the same error and performance measure. The controller should minimize
$J_{\mathrm{task}}$ while keeping the ball in the valid beam region.


\subsection{Augmented Nominal Dynamics}
We retain the hierarchical velocity-control framework of Aerial-Balance-Bench~\cite{liu2026aerialBalanceBench}.
The high-level controller regulates the ball by generating a desired vertical velocity for the drone, and the
low-level controller tracks this command. Let $c_z$ and $v_{rz}$ denote the commanded and actual drone vertical
velocities, respectively. At each sampling time $t_k$, the high-level action instantaneously updates the held
velocity command:
\begin{equation}
    u_k=\Delta c_{z,k}, \qquad
    c_z(t_k^+)=c_z(t_k^-)+u_k,
    \label{eq:incremental_velocity_command}
\end{equation}
where $t_k^-$ and $t_k^+$ denote the instants immediately before and after the update, respectively. The updated
command is held until the next sampling time.
The ball-beam dynamics are driven by $v_{rz}$ rather than directly by $c_z$. It is therefore necessary to model
the low-level actuation response between these two signals. The framework proposed in Aerial-Balance-Bench modeled this
response primarily as an action delay; here, we additionally capture its gradual first-order behavior.

We now combine the ball-beam dynamics and actuation response in one augmented model. Let $L$ and $l_r$ be the
beam and rope lengths, respectively, and let $d_b$ map the benchmark ball coordinate to the physical coordinate
along the beam. The ball has mass $m_b$, radius $r_b$, and rotational inertia $J_b$, and we define
$\gamma=m_b/(m_b+J_b/r_b^2)$. The augmented state and effective delayed command are
\begin{equation}
\begin{aligned}
    \mathbf x(t)
    &=\begin{bmatrix}
        p_b(t)&v_b(t)&\theta(t)&v_{rz}(t)&c_z(t)
    \end{bmatrix}^{\!\top},\\
    c_z^d(t)&=c_z(t-T_d),
\end{aligned}
    \label{eq:augmented_nominal_state}
\end{equation}
where $T_d$ is the response dead time. From the rope geometry, the rope angle $\beta$ and beam angular velocity $\omega$ are
\begin{equation}
\begin{aligned}
    \beta(\theta)&=\arcsin\!\left[\frac{L}{l_r}(1-\cos\theta)\right],\\
    \omega(\theta,v_{rz})&=-\frac{v_{rz}\cos\beta(\theta)}
    {L\cos\!\left(\beta(\theta)-\theta\right)}.
\end{aligned}
    \label{eq:beam_geometry}
\end{equation}
The corresponding nominal ball acceleration is
\begin{equation}
    a_b^{\mathrm{nom}}(p_b,\theta,v_{rz})
    =\gamma(p_b+d_b-L)\omega^2(\theta,v_{rz})-\gamma g\sin\theta.
    \label{eq:nominal_ball_acceleration}
\end{equation}
The low-level response is modeled by a first-order-plus-dead-time (FOPDT) relation, with time constant $\tau$,
steady-state gain $K$, and velocity bias $b$. The complete augmented nominal dynamics are
\begin{equation}
    \dot{\mathbf x}(t)
    =\mathbf f^{\mathrm{nom}}
    \!\left(\mathbf x(t),c_z^d(t)\right)
    =
    \begin{bmatrix}
        v_b(t)\\[1mm]
        a_b^{\mathrm{nom}}\!\left(p_b(t),\theta(t),v_{rz}(t)\right)\\[1mm]
        \omega\!\left(\theta(t),v_{rz}(t)\right)\\[1mm]
        \displaystyle\frac{Kc_z^d(t)+b-v_{rz}(t)}{\tau}\\[2mm]
        0
    \end{bmatrix}.
    \label{eq:augmented_nominal_dynamics}
\end{equation}
The second and third rows describe the coupled ball and beam motion, while the fourth row makes the actual drone
velocity an explicit dynamic state. The last row implements the zero-order hold described in
\eqref{eq:incremental_velocity_command}. Here, $p_b$ uses the benchmark coordinate, and $v_{rz}>0$ denotes upward
motion.

For a sampling interval $h$, $c_{z,k}^d$ is constant within one integration stage. Defining
$\bar v_{rz,k}=Kc_{z,k}^d+b$ and $a_\tau=e^{-h/\tau}$, the exact response over $s\in[0,h]$ is
\begin{equation}
\begin{aligned}
    v_{rz,k}(s)
    &=\bar v_{rz,k}+
    \left(v_{rz,k}-\bar v_{rz,k}\right)e^{-s/\tau},\\
    v_{rz,k+1}
    &=a_\tau v_{rz,k}+
    (1-a_\tau)\bar v_{rz,k}.
\end{aligned}
    \label{eq:exact_fopdt_response}
\end{equation}
The first three rows of \eqref{eq:augmented_nominal_dynamics} are integrated with RK4 using this exact response at
the start, midpoint, and end of each sampling interval. This construction defines the discrete transition maps
used below.

We identified the response model from 26 physical pseudo-random binary sequence (PRBS) trials collected at
60~Hz. We compared four structures: unit gain with delay (UGD), gain and bias with delay (GBD), first order
without delay (FO), and FOPDT. Each was fitted causally over complete response
sequences and evaluated by leave-one-trial-out validation, as shown in
Fig.~\ref{fig:velocity_response_identification}. In the upward-positive
coordinate used here, each dynamic fit was initialized with the first measured velocity and then rolled out
without injecting later velocity measurements. The reported errors therefore evaluate sequence response rather
than one-step regression. The pooled FOPDT parameters are
\begin{equation}
\begin{aligned}
    T_d &=0.0722~\mathrm{s}, & \tau&=0.1874~\mathrm{s},\\
    K   &=1.0135,           & b&=-0.00485~\mathrm{m/s}.
\end{aligned}
    \label{eq:identified_fopdt_parameters}
\end{equation}
FOPDT achieved the lowest mean held-out RMSE of $0.01659$~m/s, compared with $0.02459$~m/s for the next-best
FO model, $0.03971$~m/s for GBD, and $0.05754$~m/s for UGD. This ordering indicates that modeling the gradual
first-order response provides the main improvement over the pure-delay models, while retaining explicit dead time
yields a further gain. We therefore use FOPDT in the
augmented nominal dynamics.

\begin{figure}[t]
    \centering
    \includegraphics[width=0.5\textwidth]{figures/velocity_response_model_comparison.pdf}
    \caption{Comparison of four low-level vertical-velocity response models over 26 physical PRBS trials. Bars and error bars show means and sample standard deviations across trials, respectively.
    Training fit denotes a separate fit to each trial, while held-out trial denotes leave-one-trial-out prediction.}
    \label{fig:velocity_response_identification}
\end{figure}



%==============================================================================
\section{Structured Residual-Aware Predictive Control}
\label{sec:residual_aware_control}
\subsection{Structured Residual Dynamics}
The augmented nominal dynamics capture the main system geometry and the identified low-level actuation response,
but they cannot represent every effect in the physical ball-beam system. Small variations in the effective beam
profile, assembly, and energy dissipation can accumulate during fast and precise ball motion. In the real-world
experiments of Aerial-Balance-Bench, larger steady-state errors and degraded control performance were observed at
targets away from the beam center~\cite{liu2026aerialBalanceBench}. This target-dependent behavior indicates that
position-dependent unmodeled effects remain relevant. We therefore aim to improve control-related prediction
without claiming to recover the exact physical sources of the mismatch.

Residual augmentation has improved real-world aerial-robot control in several
settings~\cite{shi2019neural,torrente2021data,salzmann2023real}. In this system, we augment the nominal ball
acceleration with a learned residual to compensate for task-specific model mismatch. Rather than
learning an unrestricted function of the full augmented state, we impose a compact, physically structured form
that captures position-dependent and dissipative effects while remaining suitable for NMPC.

Accordingly, the residual-aware ball acceleration is
\begin{equation}
    a_b^{\mathrm{res}}(p_b,v_b,\theta,v_{rz})=a_b^{\mathrm{nom}}(p_b,\theta,v_{rz})+r_a(p_b,v_b,\theta),
    \label{eq:structured_residual_dynamics}
\end{equation}
where $r_a(p_b,v_b,\theta)$ denotes the residual acceleration term. This correction leaves the nominal beam
geometry, FOPDT velocity-response state, and command state unchanged.

We construct $r_a$ from two physically motivated components. The first accounts for position-dependent variation
in the effective local beam slope. Let $\delta\theta(p_b)$ denote the learned slope correction at ball position
$p_b$. The resulting correction to the gravity-driven ball acceleration is
\begin{equation}
    r_{\mathrm{prof}}(p_b,\theta)
    =-\gamma g\!\left[
    \sin\!\bigl(\theta+\delta\theta(p_b)\bigr)-\sin\theta\right].
    \label{eq:profile_residual}
\end{equation}
Thus, $r_{\mathrm{prof}}$ modifies the gravity projection through an effective slope change rather than adding an
arbitrary position-dependent acceleration bias.

The second component represents position-dependent dissipative effects:
\begin{equation}
    r_{\mathrm{diss}}(p_b,v_b)
    =-f_c(p_b)\tanh\!\left(\frac{v_b}{v_0}\right)-f_v(p_b)v_b.
    \label{eq:dissipative_residual}
\end{equation}
Here, $f_c(p_b)$ is a Coulomb-like acceleration magnitude in m/s$^2$, $f_v(p_b)$ is a viscous coefficient in
s$^{-1}$, and $v_0=0.025$~m/s sets the smooth transition around zero velocity. Constraining $f_c$ and $f_v$ to
be nonnegative ensures that $r_{\mathrm{diss}}$ opposes the ball motion.

Suppressing their arguments for compactness, the two components are combined and smoothly bounded as
\begin{equation}
\begin{aligned}
    r_a^{\mathrm{raw}}&=r_{\mathrm{prof}}+r_{\mathrm{diss}},\\
    r_a&=\bar r_a\tanh\!\left(\frac{r_a^{\mathrm{raw}}}{\bar r_a}\right).
\end{aligned}
    \label{eq:structured_residual_components}
\end{equation}
Here, $r_a^{\mathrm{raw}}$ is the unconstrained sum, while $\bar r_a=0.25$~m/s$^2$ sets the residual-acceleration
limit. The outer hyperbolic tangent makes the applied residual bounded and smooth for real-time differentiation.

We parameterize $\delta\theta$, $f_c$, and $f_v$ using shared normalized Gaussian radial basis functions (RBFs).
Let $M$ denote the number of bases, $\mu_i$ their centers, $\sigma_p$ their common width, and $\epsilon>0$ a
small constant in the normalization denominator. The bases and the three position-dependent quantities are
\begin{equation}
\begin{aligned}
    \phi_i(p_b)
    &=\exp\!\left[-\frac{(p_b-\mu_i)^2}{2\sigma_p^2}\right],\\
    \bar\phi_i(p_b)
    &=\frac{\phi_i(p_b)}
    {\epsilon+\sum_{j=1}^{M}\phi_j(p_b)},\\
    \delta\theta(p_b)
    &=\sum_{i=1}^{M}w_i^p\bar\phi_i(p_b),\\
    \begin{bmatrix}f_c(p_b)\\ f_v(p_b)\end{bmatrix}
    &=\sum_{i=1}^{M}\bar\phi_i(p_b)
    \begin{bmatrix}w_i^c\\ w_i^v\end{bmatrix}.
\end{aligned}
    \label{eq:normalized_position_rbf}
\end{equation}
Here, $\phi_i$ is the $i$th Gaussian basis and $\bar\phi_i$ is its normalized value. We use $M=13$, place
$\mu_i$ uniformly over the valid beam region, and set $\sigma_p=0.065$~m and
$\epsilon=10^{-12}$. The coefficients $w_i^p$, $w_i^c$, and $w_i^v$ respectively parameterize the slope
correction, Coulomb-like magnitude, and viscous coefficient. We impose
$|w_i^p|\leq\bar\delta_\theta$ with $\bar\delta_\theta=0.035$~rad,
$0\leq w_i^c\leq0.15$~m/s$^2$, and $0\leq w_i^v\leq0.50$~s$^{-1}$. These components remain structured
approximations rather than identified measurements of the beam profile or friction forces.

The residual parameters are learned offline from real closed-loop data. We form an acceleration-level training
target using the ball acceleration $\hat a_b$ obtained by differentiating a smoothed ball-position signal:
\begin{equation}
    r_a^{\mathrm{data}}(t)
    =\hat a_b(t)-a_b^{\mathrm{nom}}
    \bigl(p_b(t),\theta(t),v_{rz}(t)\bigr).
    \label{eq:residual_acceleration_target}
\end{equation}
The training objective combines four complementary terms:
\begin{equation}
    \mathcal L=
    \underset{\text{pointwise fit}}{\underline{\lambda_{\mathrm{pt}}\mathcal L_{\mathrm{pt}}}}
    +\underset{\text{integral consistency}}{\underline{\mathcal L_{\mathrm{int}}}}
    +\underset{\text{RBF smoothness}}{\underline{\lambda_{\mathrm{sm}}\mathcal L_{\mathrm{sm}}}}
    +\underset{\text{causal rollout}}{\underline{\lambda_{\mathrm{roll}}\mathcal L_{\mathrm{roll}}}}.
    \label{eq:structured_residual_training_loss}
\end{equation}
Here, $\mathcal L_{\mathrm{pt}}$ is the pointwise Huber loss between $r_a$ and the target
$r_a^{\mathrm{data}}$; $\mathcal L_{\mathrm{int}}$ enforces acceleration-integral consistency over
$H\in\{3,6,12\}$; $\mathcal L_{\mathrm{sm}}$ promotes smooth spatial variation of the RBF-parameterized
quantities; and $\mathcal L_{\mathrm{roll}}$ penalizes causal multi-step errors in predicted $p_b$ and $v_b$ over
$H\in\{6,12,30,33\}$. The $\lambda$ coefficients set the relative weights of their corresponding terms. The
first training stage sets $\lambda_{\mathrm{roll}}=0$, while the second enables the rollout term. After each
update, the coefficients are projected onto the bounds above, and the final model is selected using validation
rollout errors rather than pointwise acceleration error alone. All parameters are then frozen.


\subsection{Residual-Aware Predictive Control}
We retain the delay-aware control structure introduced in
Aerial-Balance-Bench~\cite{liu2026aerialBalanceBench}: a model-based predictor propagates the measured state
through the pending commands, and the high-level controller starts from the resulting delay-compensated state.
Our extension is to use the same augmented structured-residual dynamics in both this state predictor and the high-level NMPC controler,
as summarized in Fig.~\ref{fig:residual_aware_control_structure}. The delay queue remains confined to the state
predictor and is not modeled again over the NMPC horizon.

\begin{figure*}[t]
    \centering
    \includegraphics[width=\textwidth]{figures/residual_aware_predictive_control_structure.pdf}
    \caption{Residual-aware realization of the delay-aware control structure. The predictor advances the
    measured state through the ordered pending-command queue over the pure-delay interval, and NMPC starts from
    the resulting delay-compensated state with a time-aligned reference preview. Both the predictor and NMPC use
    the same structured residual dynamics.}
    \label{fig:residual_aware_control_structure}
\end{figure*}

\textit{1) Delay-aware state prediction:}
The model-based predictor propagates the physical state
$\boldsymbol{\xi}=[p_b,v_b,\theta,v_{rz}]^\top$ using the residual dynamics developed above. Let
$\boldsymbol{\xi}_k^m=[p_{b,k}^m,v_{b,k}^m,\theta_k^m,v_{rz,k}^m]^\top$ be its measurement, and let $D$ be the
integer delay horizon used by the controller, such that $T_d\approx Dh$. The ordered queue
$\mathcal Q_k=\{q_{0|k},\ldots,q_{D-1|k}\}$ contains the accumulated velocity commands that will become effective
over the next $D$ intervals, from oldest to newest. In particular, $q_{j|k}$ serves as the held delayed input
$c_z^d$ over $[t_{k+j},t_{k+j+1})$; it is a velocity command rather than an increment $u_k$. The predictor computes
\begin{equation}
\begin{aligned}
    \hat{\boldsymbol{\xi}}_{k|k}
    &=\boldsymbol{\xi}_k^m,\\
    \hat{\boldsymbol{\xi}}_{k+j+1|k}
    &=\mathbf G_h^{\mathrm{res}}
    \bigl(\hat{\boldsymbol{\xi}}_{k+j|k},q_{j|k}\bigr),
    \quad j=0{:}D-1,\\
    \tilde{\mathbf x}_{0|k}
    &=\begin{bmatrix}
        \hat{\boldsymbol{\xi}}_{k+D|k}^{\top}&c_z(t_k^-)
    \end{bmatrix}^{\!\top}.
\end{aligned}
    \label{eq:delay_state_prediction}
\end{equation}
Here, $\mathbf G_h^{\mathrm{res}}$ is one discrete step of the first four rows of
\eqref{eq:augmented_nominal_dynamics}, with $q_{j|k}$ substituted for $c_z^d$ and
$a_b^{\mathrm{nom}}$ replaced by $a_b^{\mathrm{res}}$ from \eqref{eq:structured_residual_dynamics}. It uses the
exact first-order velocity response and evaluates $r_a$ at every RK4 stage. The vector
$\tilde{\mathbf x}_{0|k}$ is the delay-compensated NMPC initial state: its physical component is predicted to the
end of the delay interval, whereas $c_z(t_k^-)$ is the current command accumulator before $u_k$ is applied. The
predictor estimates no new parameters and performs no online update.

Because the physical component of $\tilde{\mathbf x}_{0|k}$ corresponds to the end of the delay interval, the
NMPC reference must be aligned with the same future time:
\begin{equation}
\begin{aligned}
    p_{g,i|k}&=p_g\!\left(t_k+(D+i)h\right),\\
    v_{g,i|k}&=v_g\!\left(t_k+(D+i)h\right),
\end{aligned}
    \label{eq:delay_aligned_reference}
\end{equation}
where $i=0{:}N$ applies to both references. The reference generator supplies $v_g$, which equals $\dot p_g$
where the reference is differentiable. For
piecewise trajectories, its discrete stage value is used directly. A constant target gives
$p_{g,i|k}=p_g^\star$ and $v_{g,i|k}=0$. Thus, the raw preview must cover offsets $0$ through $D+N$. Pure delay
is not modeled again inside the high-level optimal control problem (OCP).

\textit{2) Residual-aware NMPC:}
Using the five-dimensional state in \eqref{eq:augmented_nominal_state} and the scalar command increment in
\eqref{eq:incremental_velocity_command}, the prediction model is
\begin{equation}
\begin{aligned}
    \mathbf x_{0|k}&=\tilde{\mathbf x}_{0|k},\\
    \mathbf x_{i+1|k}
    &=\mathbf F_h^{\mathrm{res}}
    \bigl(\mathbf x_{i|k},u_{i|k}\bigr),
    \quad i=0{:}N-1.
\end{aligned}
    \label{eq:residual_aware_prediction_model}
\end{equation}
At each stage, $\mathbf F_h^{\mathrm{res}}$ first applies
$c_{z,i|k}^{+}=c_{z,i|k}+u_{i|k}$ and then uses this updated command as the effective FOPDT input. The map combines
the exact first-order response, RK4 integration, and the structured residual evaluated at each RK4 stage; the
pure delay has already been consumed by the state predictor.

The nonlinear least-squares residuals for a prediction stage and the terminal stage are
\begin{equation}
    \mathbf y_{i|k}
    =\begin{bmatrix}
        p_{b,i|k}-p_{g,i|k}\\
        v_{b,i|k}-v_{g,i|k}\\
        \theta_{i|k}\\
        \omega_{i|k}\\
        v_{rz,i|k}\\
        c_{z,i|k}+u_{i|k}\\
        u_{i|k}
    \end{bmatrix},\qquad
    \mathbf y_{N|k}^{e}
    =\begin{bmatrix}
        p_{b,N|k}-p_{g,N|k}\\
        v_{b,N|k}-v_{g,N|k}\\
        \theta_{N|k}\\
        \omega_{N|k}\\
        v_{rz,N|k}\\
        c_{z,N|k}
    \end{bmatrix}.
    \label{eq:nmpc_tracking_residuals}
\end{equation}
They penalize ball tracking error, beam and drone motion, the updated velocity command, and the command increment.
No separate action-jerk term is used. Let $\mathbf s_i^l,\mathbf s_i^u\geq\mathbf 0$ denote the lower- and
upper-bound slack vectors for the soft constraints introduced below. The objective is defined as
\begin{equation}
\begin{aligned}
    J_k
    &=\frac{1}{2}\sum_{i=0}^{N-1}
    \mathbf y_{i|k}^{\top}\mathbf W\mathbf y_{i|k}+\frac{1}{2}(\mathbf y_{N|k}^{e})^{\top}
    \mathbf W_N\mathbf y_{N|k}^{e}+J_{\mathrm{slack}},\\
    J_{\mathrm{slack}}
    &=\sum_{i=0}^{N}\Bigl[
    \rho_1\mathbf 1^\top(\mathbf s_i^l+\mathbf s_i^u)+\frac{\rho_2}{2}\left(\lVert\mathbf s_i^l\rVert_2^2
    +\lVert\mathbf s_i^u\rVert_2^2\right)\Bigr],
\end{aligned}
    \label{eq:nmpc_tracking_objective}
\end{equation}
where $J_{\mathrm{slack}}$ applies linear and quadratic penalties to violations of both sides of the softened
bounds. We use the diagonal weights
\begin{equation}
\begin{aligned}
    \mathbf W&=\operatorname{diag}(20,4,0.1,0.2,0.2,0.05,0.25),\\
    \mathbf W_N&=\operatorname{diag}(100,20,0.5,1,1,0.05).
\end{aligned}
    \label{eq:nmpc_tracking_weights}
\end{equation}
with $\rho_1=\rho_2=10^4$.
Using these slacks, we soften the ball-position, beam-angle, and tracking-error bounds defined by
$\mathbf h_{i|k}=[p_{b,i|k},\theta_{i|k},p_{b,i|k}-p_{g,i|k}]^\top$. The resulting OCP is
\begin{equation}
\begin{aligned}
    \min_{\mathbf x,\mathbf u,\mathbf s^l,\mathbf s^u}\quad
    &J_k\\
    \mathrm{s.t.}\quad
    &\mathbf x_{0|k}=\tilde{\mathbf x}_{0|k},\\
    &\mathbf x_{i+1|k}=\mathbf F_h^{\mathrm{res}}
    (\mathbf x_{i|k},u_{i|k}),\\
    &-a_{\max}h\leq u_{i|k}\leq a_{\max}h,\\
    &-v_{\max}\leq
    \begin{bmatrix}v_{rz,i|k}\\c_{z,i|k}\end{bmatrix}
    \leq v_{\max},\\
    &\underline{\mathbf h}-\mathbf s_i^l
    \leq\mathbf h_{i|k}
    \leq\overline{\mathbf h}+\mathbf s_i^u,\\
    &\mathbf s_i^l,\mathbf s_i^u\geq0.
\end{aligned}
    \label{eq:residual_aware_nmpc_ocp}
\end{equation}
The dynamics and input bounds hold for $i=0{:}N-1$, the velocity bounds for $i=1{:}N$, and the softened bounds
for $i=0{:}N$. We use
$\underline{\mathbf h}=[0,-40^\circ,-0.5]^\top$ and
$\overline{\mathbf h}=[0.70,40^\circ,0.5]^\top$. The
$v_{\max}$ bounds are imposed only when a positive velocity limit is configured; the input-increment bound is
always hard.

We use $N=30$ and $h=1/60$~s. The OCP is solved with acados~\cite{verschueren2022acados} using SQP-RTI,
partial-condensing HPIPM, a Gauss-Newton Hessian approximation, and the discrete prediction map. The previous
solution is shifted for warm starting, and only $u_k=u_{0|k}^{\star}$ is applied. A nonzero solver status,
nonfinite result, or bound violation triggers $u_k=0$, which keeps the current accumulated command unchanged and
is recorded as a fallback. The issued command then updates both the command accumulator and the delay queue.
Using the same frozen residual in the predictor and NMPC keeps the two consecutive prediction intervals
dynamically consistent.



%=============================================================================
\section{Experiments}
\label{sec:experiments}

The experiments are organized around three research questions that examine unified control, control-relevant
model prediction, and real-world closed-loop performance:
\begin{itemize}
\item \textbf{RQ1:} Can nominal NMPC solve both
target-position balancing and trajectory tracking within the same framework, and how does it compare with a
non-optimization-based baseline in ideal simulation? 
\item \textbf{RQ2:} Does the structured residual
improve control-relevant prediction over the nominal model on offline real-world data? 
\item \textbf{RQ3:} Does residual-aware NMPC improve real-world target-position balancing and trajectory tracking over
nominal NMPC while retaining real-time feasibility? 
\end{itemize}

\subsection{Unified Tracking in Ideal Simulation}

\textit{1) Simulation settings:}
We conduct the simulations in Aerial-Balance-Bench~\cite{liu2026aerialBalanceBench}, which is developed using
Isaac Lab~\cite{2023_RAL_isaac_lab}. A geometric SE(3)
controller~\cite{2010_SE3_controller} tracks the drone velocity commands at 60~Hz. Each controller is evaluated
on 500 episodes of 20~s per reference type. Across
episodes, the environment samples $D_{\mathrm{sim}}\in\{6,7,8,9,10\}$,
$\tau_{\mathrm{sim}}\in[0.174,0.200]$~s, $K_{\mathrm{sim}}\in[0.95,1.05]$, and
$b_{\mathrm{sim}}\in[-0.0060,-0.0037]$~m/s. Gaussian noise is also added to the velocity response. Both
controllers use the fixed nominal values $D=8$,
$\tau=0.187$~s, $K=1.0$, and $b=-0.00485$~m/s. This setting is termed ideal because it excludes additional
ball-beam mismatch and external or position-dependent dissipative disturbances. The uncertainty is confined to the
low-level actuation response.

\textit{2) Reference design and metrics:}
We cover constant-reference balancing and dynamic tracking, with the latter represented by smooth B-spline and
transition-rich ramp-dwell trajectories:
\begin{itemize}
    \setlength{\itemsep}{0pt}
    \setlength{\parsep}{0pt}
    \setlength{\topsep}{2pt}
    \item \emph{Constant:} The goal and initial ball position are sampled independently from $[0.10,0.60]$~m,
    with at least $0.05$~m separation.
    \item \emph{Random B-spline:} A cubic 20-s trajectory uses 12 control points, starts and ends at $0.35$~m,
    and alternates its inner control points between $[0.10,0.30]$ and $[0.40,0.60]$~m to generate smooth motion.
    \item \emph{Random Ramp-dwell:} A 20-s reference starts at $0.35$~m and alternates six targets between the
    same intervals, with 4-5-s half-cycles drawn as either continuous ramps or ramp-and-dwell profiles to test
    slope changes, turning points, and dwell transitions.
\end{itemize}
The dynamic tasks initialize the ball on the reference. We evaluate the constant task using steady-state error
(SE) and convergence time (CONT). SE is the mean absolute error over the final 1~s, while CONT is the time to
enter and remain within $0.01$~m of the target, with the episode duration assigned upon nonconvergence. Dynamic
tracking is evaluated by full-episode mean absolute error (MAE) and root mean square error (RMSE).

\textit{3) Baseline:}
We implement a nonlinear feedforward-feedback (NFFB) controller as a non-optimization-based baseline. It combines
reference-acceleration feedforward and ball-state feedback, maps the desired acceleration to a beam angle through
nominal dynamic inversion, and generates the velocity command through filtering and beam-angle feedback.
For compactness, we omit time indices and evaluate the following laws using the delay-compensated
state and time-aligned reference. With $e_p=p_b-p_g$ and $e_v=v_b-v_g$, the outer loop computes
\begin{equation}
\begin{aligned}
    a_b^\star
    &=a_g-k_p e_p-k_v e_v-k_I z,\\
    \theta^\star
    &=\arcsin\!\left(
    \frac{(p_b+d_b-L)\omega^2-a_b^\star/\gamma}{g}
    \right),
\end{aligned}
    \label{eq:nffb_outer_loop}
\end{equation}
where $z$ accumulates $e_p$ with anti-windup, and $a_g$ is the limited reference acceleration
computed from consecutive preview velocities. Before inversion, $a_b^\star$ is clipped to the
acceleration range permitted by the beam-angle bounds. A second-order command filter generates
the desired beam angle $\theta_d$ and angular rate $\omega_d$:
\begin{equation}
\begin{aligned}
    \dot{\theta}_d&=\omega_d,\\
    \dot{\omega}_d
    &=\omega_f^2(\theta^\star-\theta_d)
    -2\zeta_f\omega_f\omega_d.
\end{aligned}
    \label{eq:nffb_command_filter}
\end{equation}
Here, $\omega_f$ and $\zeta_f$ specify the filter natural frequency and damping ratio.
Beam-angle feedback and the geometric velocity mapping then give
\begin{equation}
\begin{aligned}
    \omega_c&=\omega_d-k_\theta(\theta-\theta_d),\\
    c_z^\star&=-\frac{L\cos(\beta-\theta)}{\cos\beta}\,\omega_c,
\end{aligned}
    \label{eq:nffb_inner_loop}
\end{equation}
where $k_p$, $k_v$, $k_I$, and $k_\theta$ are feedback gains, and $\beta=\beta(\theta)$
follows \eqref{eq:beam_geometry}. The filtered angle and rate are bounded, and the requested
increment $c_z^\star-c_z(t_k^-)$ is limited by the command constraints to obtain $u_k$.
Similar to nominal NMPC, NFFB uses a $D=8$ state predictor based only on the augmented nominal dynamics, including the
fixed first-order response, without the residual model. NFFB is therefore delay-aware but does not solve an
online receding-horizon optimization problem. Evaluation uses identical task and response
distributions, random seed, measured states, and command constraints.

\textit{4) Results:}
The results are summarized in Table~\ref{tab:ideal_simulation_tracking}. For constant-reference balancing, the nominal NMPC reduces SE and CONT by
$60.8\%$ and $71.9\%$, respectively. For B-spline references, it reduces MAE and RMSE by $63.7\%$ and $47.7\%$,
respectively. For ramp-dwell references, the corresponding reductions are $83.5\%$ and $74.9\%$. The strongest
gains on the transition-rich task highlight the benefit of optimizing with reference preview and the predicted
actuation response. Overall, nominal NMPC handles balancing and both dynamic tasks in one formulation and
outperforms NFFB on every reported metric.

\begin{table}[t]
    \caption{Ideal-simulation tracking performance.}
    \label{tab:ideal_simulation_tracking}
    \centering
    {\footnotesize
    \setlength{\tabcolsep}{4pt}
    \renewcommand{\arraystretch}{1.25}
    \begin{tabular}{|c|c|c|c|}
        \hline
        \multirow{2}{*}{\textbf{Reference}} &
        \multirow{2}{*}{\textbf{Metric}} &
        \multicolumn{2}{c|}{\textbf{Method}} \\
        \cline{3-4}
        & & \textbf{NFFB} & \textbf{Nominal NMPC} \\
        \hline
        \multirow{2}{*}{\textbf{Constant}}
        & \textbf{SE} (cm) & $0.68{\pm}0.44$ & $\mathbf{0.27{\pm}0.16}$ \\
        \cline{2-4}
        & \textbf{CONT} (s) & $13.28{\pm}6.60$ & $\mathbf{3.73{\pm}0.79}$ \\
        \hline
        \multirow{2}{*}{\textbf{Random B-spline}}
        & \textbf{MAE} (cm) & $2.16$ & $\mathbf{0.79}$ \\
        \cline{2-4}
        & \textbf{RMSE} (cm) & $3.22$ & $\mathbf{1.68}$ \\
        \hline
        \multirow{2}{*}{\textbf{Random Ramp-dwell}}
        & \textbf{MAE} (cm) & $3.87$ & $\mathbf{0.64}$ \\
        \cline{2-4}
        & \textbf{RMSE} (cm) & $4.51$ & $\mathbf{1.13}$ \\
        \hline
    \end{tabular}}
\end{table}



\subsection{Offline Residual-Dynamics Prediction}

\textit{1) Dataset:}
We use real-world closed-loop data processed as control snapshots at approximately
60~Hz. The dataset is split by complete trials into 79 for training, 6 for validation, and 6 for held-out testing;
each split covers constant, sinusoidal, triangular, and trapezoidal references. This trial-level division prevents
windows from the same trial from appearing in both model fitting and testing. All residual parameters are selected
using only the training and validation data, with the held-out test trials reserved for final evaluation.

\textit{2) Causal rollout protocol:}
We emphasize multi-step ball-position and ball-velocity prediction because NMPC optimizes these state trajectories
over a finite horizon, whereas one-step acceleration error alone does not reveal how model errors accumulate after
integration. The 6 held-out trials yield 2,028 prediction windows shared by the nominal and structured-residual
models. Each window initializes both rollouts from the same measured state and pending-command queue, after which
they use the same recorded command sequence, FOPDT response, geometric parameters, and RK4 integration. Neither
rollout is reset using future state measurements. Instead, each model is propagated once for 60 steps, with
$r_a$ in the structured model evaluated from the predicted $p_b$, $v_b$, and $\theta$ at every RK4 stage. We evaluate the
endpoint position and velocity errors at
$H\in\{6,12,18,24,30,36,42,48,54,60\}$ steps, corresponding to $0.1$--$1.0$~s at 60~Hz in $0.1$-s intervals.
In particular, $H=30$ matches the NMPC horizon, while $H=60$ measures error accumulation over one second.

\textit{3) Results:}
Fig.~\ref{fig:offline_residual_rollout_rmse} shows that the structured-residual model achieves lower position
and velocity RMSE at every evaluated horizon. At $H=30$, it reduces position and velocity RMSE by $9.3\%$ and
$12.0\%$, respectively. At $H=60$, the reductions increase to $19.7\%$ and $20.8\%$. The growing margin over
longer rollouts indicates that the structured residual slows the accumulation of state-prediction error over
control-relevant horizons. These held-out causal-rollout results answer RQ2 positively, while the corresponding
closed-loop control benefit is evaluated in the next subsection.

\begin{figure}[t]
    \centering
    \includegraphics[width=\columnwidth]{figures/offline_residual_rollout_rmse.pdf}
    \caption{Held-out causal rollout prediction over 2,028 common windows. Both models use the same measured
    initial state, pending-command queue, and subsequent recorded commands. The left and right panels show
    ball-position and ball-velocity RMSE, respectively; the horizon is expressed in 60-Hz control steps.}
    \label{fig:offline_residual_rollout_rmse}
\end{figure}



\subsection{Real-World Closed-Loop Evaluation}

\textit{1) Experimental setup:}
We evaluate the nominal and residual-aware NMPC controllers on the physical platform shown in
Fig.~\ref{fig:real_world_system}. The drone is equipped with an NVIDIA Jetson Orin NX onboard computer and a
Pixhawk 6C Mini flight controller running PX4~\cite{2015_px4} for low-level velocity control.
A rope connects the drone to one end of the beam, while the opposite end
is supported at a stationary pivot. Vertical drone motion changes the beam inclination and guides the freely
rolling ball, whose position is measured by a laser distance sensor.
The high-level controllers run at 60~Hz and use the same state,
objective, constraints, $N=30$, $D=3$, identified FOPDT parameters, and deployment interface. Their only
prediction-model difference is that residual-aware NMPC enables the structured residual in both
the state predictor and NMPC. We test three constant targets,
$p_g^\star\in\{0.15,0.35,0.55\}$~m, and sine, triangle, and trapezoid references with periods of 6 and 8~s.
The constant targets cover the two off-center regions and the beam center, while the dynamic set varies both
motion speed and the sharpness of reference transitions. Residual parameters remain fixed throughout the tests.
Each of the nine conditions is repeated three times for each controller.

\begin{figure}[t]
    \centering
    \includegraphics[width=0.8\columnwidth,height=0.5\columnwidth]{figures/real_world_system.png}
    \caption{Physical drone-ball platform. The drone carries an NVIDIA Jetson Orin NX and a Pixhawk 6C Mini
    and is connected to the beam by a rope. A laser distance sensor measures the ball position along the beam.}
    \label{fig:real_world_system}
\end{figure}

\textit{2) Evaluation metrics:}
Each trial lasts 30~s. For constant targets, we retain the SE and CONT definitions used in the ideal-simulation
evaluation, but use a 5-s terminal window for SE and a $0.03$-m convergence threshold for CONT; a nonconverged
trial is assigned 30~s. Dynamic tracking is evaluated over
the complete trial using position MAE, RMSE, and response gain $G=A_b/A_g$, where $A_b$ and $A_g$ are the fitted
ball-response and reference amplitudes. For an error metric $E$ and for gain, respectively, the reported
improvements are
\begin{equation}
\begin{aligned}
    \mathcal{I}_E
    &=100\frac{E_{\mathrm{nom}}-E_{\mathrm{res}}}{E_{\mathrm{nom}}}\%,\\
    \mathcal{I}_G
    &=100\frac{|1-G_{\mathrm{nom}}|-|1-G_{\mathrm{res}}|}
    {|1-G_{\mathrm{nom}}|}\%.
\end{aligned}
\label{eq:real_world_improvement}
\end{equation}
Each entry is the mean of three trials; positive improvement denotes better residual-aware performance.

\textit{3) Tracking results:}
Fig.~\ref{fig:real_world_closed_loop_tracking} gives representative traces, while
Tables~\ref{tab:real_world_constant_tracking} and~\ref{tab:real_world_dynamic_tracking} aggregate all trials.
In both tables, Nom. and Res. denote nominal and residual-aware NMPC, respectively.
For constant targets, residual-aware NMPC reduces overall SE by $55.2\%$ and CONT by $46.7\%$. The largest SE reduction, $77.6\%$, occurs at $0.55$~m, whereas
the $0.35$-m SE changes by only $-0.8\%$. Thus, the structured residual mainly corrects the larger off-center
bias without noticeably changing the already accurate center-target result. Across the six dynamic conditions,
the overall MAE and RMSE decrease by $31.1\%$ and $16.5\%$. The response gain increases from $0.855$ to $0.950$,
which reduces the amplitude deviation from unity by $65.2\%$. Every dynamic condition improves in MAE and gain
deviation, and five of six improve in RMSE. The exception is the 6-s trapezoid, whose RMSE increases by $3.3\%$;
the traces and error metrics suggest that local errors around its sharp transitions remain challenging.
The representative traces show the same main pattern: the center-directed offsets at the two outer constant
targets are reduced, and the dynamic peaks and valleys follow the requested amplitude more closely.

\begin{figure*}[t]
    \centering
    \includegraphics[width=\textwidth]{figures/real_world_closed_loop_tracking.pdf}
    \caption{Representative real-world tracking trials. Constant-reference trials are selected by the median SE
    and dynamic trials by the median MAE for each controller and condition. The dynamic panels use an 8-s
    period.}
    \label{fig:real_world_closed_loop_tracking}
\end{figure*}

\begin{table}[t]
    \caption{Real-world constant-target results.}
    \label{tab:real_world_constant_tracking}
    \centering
    {\renewcommand{\arraystretch}{1.2}
    \begin{tabular}{|c|ccc|ccc|}
        \hline
        \multirow{2}{*}{\textbf{Target}} &
        \multicolumn{3}{c|}{\textbf{SE} (cm)} &
        \multicolumn{3}{c|}{\textbf{CONT} (s)} \\
        \cline{2-7}
        & \textbf{Nom.} & \textbf{Res.} & \textbf{Imp.} (\%)
        & \textbf{Nom.} & \textbf{Res.} & \textbf{Imp.} (\%) \\
        \hline
        $\mathbf{0.15}$~\textbf{m} & 3.035 & \textbf{2.404} & \textbf{$+20.8$}
        & 21.108 & \textbf{12.810} & \textbf{$+39.3$} \\
        $\mathbf{0.35}$~\textbf{m} & \textbf{0.915} & 0.921 & $-0.8$
        & 7.510 & \textbf{6.712} & \textbf{$+10.6$} \\
        $\mathbf{0.55}$~\textbf{m} & 6.954 & \textbf{1.555} & \textbf{$+77.6$}
        & 30.000 & \textbf{11.700} & \textbf{$+61.0$} \\
        \hline
        \textbf{Overall} & 3.634 & \textbf{1.627} & \textbf{$+55.2$}
        & 19.539 & \textbf{10.407} & \textbf{$+46.7$} \\
        \hline
    \end{tabular}}
\end{table}

\begin{table*}[t]
    \caption{Real-world dynamic-tracking results.}
    \label{tab:real_world_dynamic_tracking}
    \centering
    {\renewcommand{\arraystretch}{1.2}
    \begin{tabular}{|c|ccc|ccc|ccc|}
        \hline
        \multirow{2}{*}{\textbf{Reference}} &
        \multicolumn{3}{c|}{\textbf{MAE} (cm)} &
        \multicolumn{3}{c|}{\textbf{RMSE} (cm)} &
        \multicolumn{3}{c|}{\textbf{Response gain}} \\
        \cline{2-10}
        & \textbf{Nom.} & \textbf{Res.} & \textbf{Imp.} (\%)
        & \textbf{Nom.} & \textbf{Res.} & \textbf{Imp.} (\%)
        & \textbf{Nom.} & \textbf{Res.} & \textbf{Error red.} (\%) \\
        \hline
        \textbf{Sine. 6~s} & 2.043 & \textbf{1.236} & \textbf{$+39.5$}
        & 2.893 & \textbf{2.244} & \textbf{$+22.4$}
        & 0.869 & \textbf{0.969} & \textbf{$+76.2$} \\
        \textbf{Sine. 8~s} & 2.293 & \textbf{1.402} & \textbf{$+38.8$}
        & 2.802 & \textbf{2.287} & \textbf{$+18.4$}
        & 0.850 & \textbf{0.955} & \textbf{$+70.3$} \\
        \textbf{Triangle. 6~s} & 2.198 & \textbf{1.648} & \textbf{$+25.0$}
        & 2.806 & \textbf{2.658} & \textbf{$+5.3$}
        & 0.849 & \textbf{0.925} & \textbf{$+50.5$} \\
        \textbf{Triangle. 8~s} & 2.087 & \textbf{1.365} & \textbf{$+34.6$}
        & 2.530 & \textbf{1.941} & \textbf{$+23.3$}
        & 0.841 & \textbf{0.938} & \textbf{$+60.8$} \\
        \textbf{Trapezoid. 6~s} & 2.552 & \textbf{2.350} & \textbf{$+7.9$}
        & \textbf{3.635} & 3.756 & $-3.3$
        & 0.874 & \textbf{0.947} & \textbf{$+58.0$} \\
        \textbf{Trapezoid. 8~s} & 2.715 & \textbf{1.563} & \textbf{$+42.4$}
        & 3.486 & \textbf{2.276} & \textbf{$+34.7$}
        & 0.851 & \textbf{0.965} & \textbf{$+76.3$} \\
        \hline
        \textbf{Overall} & 2.315 & \textbf{1.594} & \textbf{$+31.1$}
        & 3.025 & \textbf{2.527} & \textbf{$+16.5$}
        & 0.855 & \textbf{0.950} & \textbf{$+65.2$} \\
        \hline
    \end{tabular}}
\end{table*}

\textit{4) Real-time feasibility:}
Across the 27 trials for each controller, the mean/99th-percentile/maximum total computation times are
$1.583/2.704/3.452$~ms for nominal NMPC and $1.746/2.972/3.731$~ms for residual-aware NMPC. Both are well below
the $16.67$-ms control period. All 48,616 residual-aware updates return zero solver status, with no fallback,
soft deadline miss, or hard deadline miss; the nominal controller likewise records none. Therefore, under the
tested platform, references, and independent repeats, residual-aware NMPC improves overall balancing and dynamic
tracking accuracy, reduces amplitude attenuation, and retains 60-Hz feasibility. The isolated 6-s trapezoid RMSE
increase also shows that this conclusion does not imply an improvement in every individual metric.



%=============================================================================
\section{CONCLUSIONS}
\label{sec:conclusions}
This paper developed a unified residual-aware predictive-control framework for indirect drone--ball
manipulation. Target-position balancing is treated as constant-reference tracking, so one NMPC formulation
handles static targets and time-varying trajectories. Its augmented nominal model combines ball--beam mechanics
with an identified FOPDT low-level actuation response, while a pending-command predictor compensates for pure
delay. Remaining mismatch is represented by a bounded and smooth ball-acceleration residual with a
position-dependent beam-profile correction and a friction-inspired dissipative term. The same frozen residual is
used in the predictor and NMPC without changing the identified command channel.

The experiments provide three levels of evidence. In ideal simulation, nominal NMPC achieved lower error than
NFFB on all reported balancing and tracking metrics. On held-out physical data, the structured residual reduced
one-second position and velocity rollout RMSE by $19.7\%$ and $20.8\%$. In physical closed-loop experiments,
residual-aware NMPC reduced overall constant-target SE by $55.2\%$ and dynamic MAE by $31.1\%$, while raising
response gain from $0.855$ to $0.950$. It was deployed at 60~Hz. Under residual-aware NMPC, the 6-s trapezoid
RMSE increased by $3.3\%$, so the observed benefit does not extend to every individual metric. Taken together,
these results support
structured, control-oriented residual dynamics for improving multi-step prediction and unified closed-loop
tracking on this platform.

The study remains limited to the vertically constrained Aerial-Balance-Bench system, the tested reference set,
and a residual learned offline from a fixed dataset. The learned components are structured approximations rather
than direct measurements of the beam profile or friction. Future work will study safe online updates, wider
operating conditions, and indirect aerial manipulation with less constrained vehicle and support motion.







\bibliographystyle{ieeetr}
\bibliography{references}

\end{document}
